
\chapter{Introducción}
%% Este 'capitulo' o sección del documento no  numera ni secciones o subsecciones, se utiliza el esquema de libro para ordenar la información

El cuidado de las gallinas, al igual que el de cualquier otro animal, es una labor que consume muchos recursos y tiempo. Hoy en día, existe gente que ha comenzado a encargarse del cuidado de estos animales, montando sistemas de producción de traspatio para consumo personal o comercialización del huevo, los cuales exigen su debido cuidado.

Las gallinas Hy-Line Brown son aves de postura que se promocionan con un alto potencial genético, el cual se puede aprovechar con los cuidados adecuados que se describen en su guía de manejo, de acuerdo con la raza\cite{Hy-Line} .

En el presente trabajo propone el desarrollo de un sistema automático de alimentación de gallinas ponedoras que permita optimizar la gestión y reducir los costos asociados a esta actividad, junto con un sistema de monitoreo y registro de postura.

Este sistema está diseñado como un sistema automático de alimentación, en el cual se pueda confiar para la realización de la actividad. Además, al ser una actividad económica, registrará la postura de las gallinas y procesará dicha información para generar reportes que permitan una toma de decisiones rápida y eficiente ante cualquier alteración en la producción.

\newpage
\section{Definición del problema}

Actualmente se ha comenzado a buscar que la producción del huevo y el trato a las gallinas sea de manera digna a través de la producción de huevo de libre pastoreo. En este modelo se busca un trato más natural a las gallinas, teniéndolas en un espacio amplio que les permita vivir adecuadamente.

Sin embargo, el cuidado de estos seres implica una gran labor, que para una persona sería una actividad muy desgastante a mayor volumen de aves. Actividades como lo son la alimentación de las gallinas, la evaluación del agua que consumen y la recolección del huevo son repetitivas a través de un día. Por lo que es debida la implementación de un sistema que permita al usuario agilizar los procesos sin ser invasivos.

Es importante mencionar que, a pesar de ser un proceso que busca el bienestar de los animales, no deja de ser una actividad económica, en la que se busca la mayor producción evitando pérdidas; es por ello que un problema es identificar a aquellas gallinas improductivas y tomar una decisión sobre un posible descarte en la producción de huevo o, en su defecto, obtener una causa y dar solución de ser posible.

Existen sistemas de alimentación diseñados para una cantidad más grande de gallinas, como los producidos por Hotraco Agri\cite{b-2024} o Poultry Plast \cite{veterinaria-digital-sa-2021} \cite{unknown-author-2021}. Los cuales son pensados para galpones tecnificados con más de 1000 gallinas, haciéndolos rentables en casos de alta producción; no existiendo versiones aptas para sistemas semi-tecnificados con menos de 100 gallinas.

Es debido diseñar un sistema que sea pensado en baja producción de huevo, que facilite la actividad de la alimentación y lo haga de manera eficiente, evitando el desperdicio de alimento y permita a las personas encargadas del cuidado de los animales enfocarse en otras actividades que favorezcan el rendimiento.

\newpage

\section{Justificación}

La producción de traspatio es una actividad que implica atención y tiempo. Se establece en el manual que, para lograr una buena producción, es debido a un ambiente adecuado y atención ante la alimentación, hidratación y nutrición de los animales. De acuerdo con la Organización Mundial de Sanidad Animal (OMSA) \cite{world-organisation-for-animal-health-2024}, los animales terrestres requieren de Bienestar Animal; para lograrlo, es necesario cumplir con los siguientes cinco parámetros:

\begin{itemize}
    \item Libre de hambre, sed y desnutrición.
    \item Libre de temor y angustia.
    \item Libre de molestias físicas y térmicas.
    \item Libre de dolor, de lesión y enfermedad.
    \item Libre de manifestar un comportamiento natural.
\end{itemize}

En la Ciudad de México se propuso la idea de automatizar un sistema de producción de traspatio el cual actualmente cuenta con 31 gallinas, dicho sistema será atendido por un número reducido de personas. Considerando el limitado tiempo que expresan tener, se solicita un sistema que pueda apoyar en el manejo y atención de los animales, cumpliendo con las directrices de Bienestar Animal.

Se plantea un sistema que automatice la tarea de alimentación que proporcione el alimento suficiente y adecuado para cada gallina; y en vista de que se trata de una actividad pensada para la venta del huevo, se propone un sistema de monitoreo de producción de este que pueda identificar a aquellas gallinas que no estén produciendo en el tiempo adecuado.

El sistema de alimentación está pensado como una adaptación de los alimentadores tecnificados, que están diseñados para granjas con una población mucho mayor de gallinas. Esto, para que el personal encargado del cuidado pueda enfocarse en otras cuestiones del proceso, haciendo más eficaz y eficiente la tarea de alimentación.

Es importante mencionar que el objetivo del proyecto es una propuesta que se apoya mucho en los trabajos ya existentes; sin embargo, varios están enfocados a granjas industriales que manejan una cantidad mucho mayor a la planteada. Lo que podría crear un producto más alcanzable para aquellos productores de huevo que lo hacen desde su casa con sistemas de producción de traspatio o semi-tecnificados, buscando una posterior comercialización de un sistema similar. Para ello, se ha contactado con Médicos Veterinarios Zootecnistas que apoyen en que el sistema cumpla con las directrices.

\newpage
\section{Objetivo}
%\section{Objetivos}

\subsection{Objetivos Trabajo Terminal 1}
\begin{enumerate}
    \item Diseñar con software CAD el sistema estructural contenedor y el sistema de movimiento que distribuirá el alimento.
    \item Diseñar los sistemas eléctricos y electrónicos para pesar y transportar el alimento.
    \item Diseñar el sistema de monitoreo de la postura de huevo, que proporcione la información adecuada identificando a cada gallina y la frecuencia con la que pone huevo.
    \item Diseñar una interfaz hombre máquina que permita ingresar la cantidad de alimento a distribuir y la frecuencia con la que se debe repartir.
    \item Validar computacionalmente los sistemas integrados con ayuda de software CAD.
\end{enumerate}

PROPUESTA OBJETIVOS TRABAJO TERMINAL 2
1.	Construir el sistema estructural y el sistema de movimiento.
2.	Ensamblar el sistema estructural y el sistema de movimiento.
3.	Construir el sistema de monitoreo.
4.	Validar circuitos electrónicos para cada uno de los sistemas.
5.	Implementar la interfaz hombre máquina para la distribución de alimento.
6.	Integrar el sistema de alimentación y de monitoreo de la postura de huevo.
7.	Realizar pruebas de movimiento, distribución y monitoreo.


\subsection{Objetivo general}
Diseñar e implementar, en un sistema de traspatio, un sistema automático de nutrición y de monitorización de la producción de huevo.

\subsection{Objetivos específicos}

\subsubsection{Trabajo Terminal I}

\begin{itemize}
    \item Diseñar un sistema mecatrónico que integre de manera eficiente los componentes mecánicos, electrónicos y de control para la entrega precisa de alimentos y el monitoreo.
    \item Desarrollar un sistema que utilice sensores y algoritmos de control para ajustar automáticamente la cantidad de alimento entregada según las necesidades de las gallinas.
    \item Crear componentes y subsistemas modulares que permitan una fácil expansión o adaptación del sistema en el futuro.
    \item Garantizar que los componentes mecánicos sean robustos y duraderos para resistir el uso continuo en un entorno agrícola.
    \item Diseñar el sistema eléctrico y mecánico de manera que minimice el consumo de energía y maximice la eficiencia operativa.
    \item Diseñar el sistema de manera que permita un mantenimiento sencillo y, en caso de fallas, facilitar las labores de reparación.

  \end{itemize}

\subsubsection{Trabajo Terminal II}

\begin{itemize}
    \item Fabricar y ensamblar los componentes mecánicos, como las tolvas móviles y los sistemas de distribución de alimentos, junto con la electrónica necesaria para su control.
    \item Programar el firmware y el software que controlarán los actuadores y sensores del sistema, permitiendo una operación precisa y automática.
    \item Implementar sensores avanzados, como cámaras de visión por computadora, para identificar gallinas no productivas y recopilar datos de alimentación y comportamiento.
    \item Asegurar que el sistema sea capaz de realizar movimientos precisos y suaves para la entrega de alimentos y el control del ambiente.
    \item Establecer una comunicación efectiva entre los componentes electrónicos y permitir la supervisión remota a través de redes y dispositivos electrónicos.
    \item Realizar pruebas exhaustivas para verificar la correcta interacción entre los componentes mecánicos y electrónicos, así como su desempeño en condiciones reales.
\end{itemize}

%\section{Enfoque mecatrónico}



%%%%%%%%%%%%%%%%%%%%%
\newpage
\section{Antecedentes}

% Please add the following required packages to your document preamble:
% \usepackage{multirow}
% \usepackage{graphicx}
\begin{table}[htbp]
\centering
\caption{Antecedentes}
\label{tab:Ante}
\resizebox{\textwidth}{!}{%
\begin{tabular}{|c|c|c|c|}
\hline
Nombre &
  Características &
  País &
  \begin{tabular}[c]{@{}c@{}}Escuela / \\ Empresa\end{tabular} \\ \hline
\multirow{2}{*}{\begin{tabular}[c]{@{}c@{}}Análisis y construcción de un sistema automatizado \\ para la distribución de alimento para gallinas.\end{tabular}} &
  \begin{tabular}[c]{@{}c@{}}-Distribución de alimento \\ por tolvas móviles\end{tabular} &
  \multirow{2}{*}{México} &
  \multirow{2}{*}{\begin{tabular}[c]{@{}c@{}}Tésis, \\ UPIITA-IPN\end{tabular}} \\ \cline{2-2}
 &
  \begin{tabular}[c]{@{}c@{}}-Sistema tecnificado\end{tabular} &
   &
   \\ \hline
\multirow{3}{*}{\begin{tabular}[c]{@{}c@{}}Alimentador automático\\ para perros\end{tabular}} &
  \begin{tabular}[c]{@{}c@{}}-Almacenamiento \\ de alimento en tolva\end{tabular} &
  \multirow{3}{*}{México} &
  \multirow{3}{*}{\begin{tabular}[c]{@{}c@{}}Tésis, \\ UPIITA-IPN\end{tabular}} \\ \cline{2-2}
 &
  \begin{tabular}[c]{@{}c@{}}-Distribución de alimento\end{tabular} &
   &
   \\ \cline{2-2}
 &
  \begin{tabular}[c]{@{}c@{}}-Pesaje de alimento \\ por báscula\end{tabular} &
   &
   \\ \hline
\multirow{3}{*}{\begin{tabular}[c]{@{}c@{}}Control de \\ alimento y agua.\end{tabular}} &
  \begin{tabular}[c]{@{}c@{}}-Pesaje de alimento \\ por báscula\end{tabular} &
  \multirow{3}{*}{Países Bajos} &
  \multirow{3}{*}{\begin{tabular}[c]{@{}c@{}}Trabajo de \\ empresa, \\ Hotraco Agri\end{tabular}} \\ \cline{2-2}
 &
  \begin{tabular}[c]{@{}c@{}}-Sistema tecnificado\end{tabular} &
   &
   \\ \cline{2-2}
 &
  \begin{tabular}[c]{@{}c@{}}-Distribución de alimento por \\ recipientes de alimentación.\end{tabular} &
   &
   \\ \hline
\begin{tabular}[c]{@{}c@{}}Función de Producción \\ en gallina de postura\\ Lohmann Brown\end{tabular} &
  \begin{tabular}[c]{@{}c@{}}-Estudio sobre la producción \\ de huevo en una \\ granja avícola\end{tabular} &
  México &
  Tésis, UAEM \\ \hline
\multirow{3}{*}{\begin{tabular}[c]{@{}c@{}}Automatización \\ de planta avícola.\end{tabular}} &
  \begin{tabular}[c]{@{}c@{}}-Distribución de alimento \\ por recipientes de alimentación.\end{tabular} &
  \multirow{3}{*}{Colombia} &
  \multirow{3}{*}{\begin{tabular}[c]{@{}c@{}}Tésis,\\ Universidad Distrital \\ Francisco José De Caldas\end{tabular}} \\ \cline{2-2}
 &
  \begin{tabular}[c]{@{}c@{}}-Sistema de traspatio\end{tabular} &
   &
   \\ \cline{2-2}
 &
  \begin{tabular}[c]{@{}c@{}}- Control de temperatura\end{tabular} &
   &
   \\ \hline
\multirow{3}{*}{\begin{tabular}[c]{@{}c@{}}Diseño de los sistemas de automatización \\ para la ampliación de una granja avícola.\end{tabular}} &
  \begin{tabular}[c]{@{}c@{}}-Sistema tecnificado\end{tabular} &
  \multirow{3}{*}{Ecuador} &
  \multirow{3}{*}{\begin{tabular}[c]{@{}c@{}}Tésis, \\ Escuela Politécnica \\ Nacional\end{tabular}} \\ \cline{2-2}
 &
  \begin{tabular}[c]{@{}c@{}}-Distribución de alimento\end{tabular} &
   &
   \\ \cline{2-2}
 &
  \begin{tabular}[c]{@{}c@{}}-Pesaje de alimento \\ por báscula\end{tabular} &
   &
   \\ \hline
\end{tabular}%
}
\end{table}



%%%%%%fin del archivo
\endinput 